\chapter*{Remerciements}\addcontentsline{toc}{chapter}{Remerciements}


Au terme de ce stage, il m’est particulièrement agréable d’exprimer ma profonde gratitude à toutes celles et à tous ceux qui, de près ou de loin, ont contribué à la bonne réalisation de ce rapport.
Je tiens à exprimer mes sincères remerciements au Dr Abdou \textsc{DIOUF}, Directeur Général de l’ANSD, ainsi qu’au Dr Momath \textsc{CISSÉ},Directeur Général Adjoint de l’ANSD et à M. Fodé \textsc{DIÈME}, directeur de la DMCI pour m’avoir fait l’honneur de m’accueillir au sein de cette institution.\\
J’adresse également mes remerciements à M. Moussa \textsc{FALL}, Chef de la division de la Méthodologie et de l’Innovation, ainsi qu’à M. \textsc{NDIASSE}, Responsable du Bureau des Méthodes et Normes (BMN), structure au sein de laquelle j’ai eu le privilège d’effectuer mon stage.
Mes remerciements les plus chaleureux vont également à mon encadreur, M. Mamadou \textsc{GUEYE} ainsi qu'à M. Ouseynou \textsc{MBAYE}, M. Cheikh \textsc{GUEYE} Ingénieurs statisticiens économistes à la DMCI, pour leur disponibilité, leurs conseils avisés et le suivi constant dont j’ai bénéficié tout au long de ces deux mois.\\  
Je souhaite aussi exprimer ma reconnaissance à l’ensemble du personnel de la DMCI pour leur accueil chaleureux, leurs échanges enrichissants et leur ouverture d’esprit.\\
Je ne saurais oublier de remercier l’ensemble du corps administratif et professoral de l’ENSAE, en particulier M. Idrissa \textsc{DIAGNE}, Directeur de l’ENSAE, Dr Alassane \textsc{AW}, Responsable de la filière des Analystes Statisticiens (AS), Dr. Mamadou \textsc{BALDÉ}, Dircteur des études et stages.\\
Enfin, j’adresse mes remerciements les plus chaleureux à mes camarades de la promotion \textsc{AS 2023--2026}, pour l’esprit de solidarité et de fraternité qui nous unit, ainsi qu’à toutes les personnes qui, bien que non citées nommément, ont apporté leur soutien et contribué, d’une manière ou d’une autre, à l’aboutissement de ce travail.  



\chapter*{Avant-propos}\addcontentsline{toc}{chapter}{Avant-propos}

Située au Sénégal et précisément à Dakar, l’École nationale de la Statistique et de
l’Analyse Économique Pierre Ndiaye (ENSAE) est un établissement d’enseignement
supérieur rattaché à l’Agence nationale de la Statistique et de la Démographie
(ANSD). Elle est en relation avec trois autres écoles du même statut à savoir
l’École nationale Supérieure de la Statistique et de l’Économie appliquée (ENSEA)
d’Abidjan l’Institut Sous régional de la Statistique et de l’Économie Appliquée
(ISSEA) de Yaoundé et l’Ecole nationale d’Economie appliquée et de Management
(ENEAM) de Cotonou avec lesquelles elle forme le Réseau des Ecoles de Statistiques
Africaines (RESA).
L’offre de formation initiale de l’ENSAE comprend la formation des Ingénieurs
statisticiens économistes (ISE) en cycles court et long et des Analystes statisticiens (AS). L’admission s’y fait par voie de concours. Elle dispose également des Masters spécialisés tels que le master « Aide à la Décision et Evaluation des Politiques Publiques » (ADEPP) et le Master Agricole.
Afin de garantir une formation de qualité, il est prévu que les élèves effectuent au moins un stage durant leur cursus. Ainsi, les élèves analystes statisticiens et Ingénieurs statisticiens économistes cycle long réalisent un stage d’immersion à la fin de leur deuxième année d’études. Ce stage vise à les familiariser avec le milieu professionnel et leur offrir l’opportunité de mettre en pratique les connaissances théoriques acquises. Il aboutit à la rédaction d’un rapport de stage sur un thème 
spécifique, qui sera ensuite présenté lors d'une soutenance devant un jury.
Le présent rapport, intitulé « Construction d'un indicateur ODD national au Sénégal » s’inscrit dans ce cadre et a été rédigé à l’issue d’un stage de deux mois à l’Agence nationale de la Statistique et de la Démographie au sein de la Direction de la Méthodologie, de la Coopération statistique et de l'Innovation.

